El desarrollo de SocratiCode se ha llevado a cabo siguiendo una metodología de trabajo basada en los principios del desarrollo ágil, adaptada a las particularidades de un proyecto realizado por un único desarrollador. En lugar de optar por un modelo tradicional en cascada (\textit{waterfall}), se ha priorizado un enfoque iterativo e incremental que permite una mayor flexibilidad ante los cambios de requisitos y una validación constante del progreso \cite{larman2003iterative}.

En una primera fase, el objetivo principal fue el diseño y construcción de un \textbf{producto mínimo viable (MVP)} \cite{hokkanen2016mvp}. Este MVP consistía en la integración básica de la interfaz de chat con el modelo de lenguaje ejecutado localmente, sin la complejidad añadida de la gestión de usuarios, la persistencia en base de datos o el aislamiento del motor de ejecución. La obtención temprana de este prototipo funcional fue crucial, ya que permitió validar la viabilidad técnica del \textit{streaming} de respuestas y la capacidad del modelo para seguir un \textit{prompt} socrático inicial.

A partir de este MVP, el ciclo de vida del proyecto se estructuró en iteraciones cortas o \textit{sprints}. Cada \textit{sprint} comprendía un ciclo completo de análisis de requisitos, diseño técnico, implementación y pruebas de las funcionalidades seleccionadas, como la implementación del entorno seguro (\textit{sandbox}), la creación del panel de configuración de modelos, o la introducción del flujo asíncrono de moderación. Al finalizar cada ciclo, el sistema se encontraba en un estado estable y desplegable.

Para garantizar la calidad técnica y la adecuación a los objetivos pedagógicos, se establecieron reuniones periódicas de seguimiento con el tutor del proyecto. En dichas reuniones, se demostraban los avances del \textit{sprint} finalizado, se analizaban los bloqueos técnicos encontrados y, fundamentalmente, se validaban las decisiones arquitectónicas antes de avanzar a la siguiente fase. Esta retroalimentación constante ha sido indispensable para afinar el comportamiento del tutor socrático y asegurar que el sistema evolucionara alineado con las expectativas académicas y de seguridad planteadas.